% from ega3-1-fr.tex
\oldpage[III]{138}

\begin{corollary}[4.4.11]
\label{III.4.4.11-fr}
Soient $Y$ un préschéma localement noethérien, $f:X\to Y$ un morphisme
propre. L'ensemble $U$ des points $y\in Y$ tels que $f^{-1}(y)$ soit discret
est ouvert dans $Y$, et le morphisme $f^{-1}(U)\to U$ restriction de $f$ est
fini. En particulier, un morphisme propre et quasi-fini $X\to Y$ est fini.
\end{corollary}

\begin{proof}
En effet, le complémentaire de $U$ dans $Y$ est l'image par $f$ de $X-X'$ qui
est fermé dans $X$ en vertu de (\hyperref[III.4.4.10-fr]{4.4.10})~; comme $f$
est une application fermée, $U$ est ouvert. En outre, il résulte de
(II, 6.2.2) que $f^{-1}(y)$ est fini pour tout $y\in U$~; comme le morphisme
$f^{-1}(U)\to U$, restriction de $f$ est propre (II, 5.4.1), il est fini en
vertu de (\hyperref[III.4.4.2-fr]{4.4.2}).
\end{proof}

\begin{env}[4.4.12]
\label{III.4.4.12-fr}
\emph{Remarques} (4.4.12). ---
\begin{enumerate}
\item[(i)] Comme on l'a annoncé dans (II, 6.2.7), nous montrerons au chap. V
que si $Y$ est localement noethérien, tout morphisme quasi-fini et séparé
$f:X\to Y$ est quasi-affine, donc quasi-projectif. Il s'ensuivra que, dans le
Main Theorem (\hyperref[III.4.4.3-fr]{4.4.3}) la conclusion reste valable
lorsqu'on suppose seulement $f$ séparé et de type fini. En effet, il résulte
de (\hyperref[III.4.4.10-fr]{4.4.10}) que $X'$ est ouvert dans $X$, et comme
$X$ est localement noethérien, la restriction de $f$ à $X'$ est encore de
type fini (I, 6.3.5), donc quasi-fini par définition de $X'$, et évidemment
séparé~; on peut donc appliquer (\hyperref[III.4.4.3-fr]{4.4.3}) à cette
restriction, d'où la conclusion.

\item[(ii)] Nous donnerons au chap. IV une démonstration plus élémentaire de
(\hyperref[III.4.4.10-fr]{4.4.10}), utilisant la théorie de la dimension.
\end{enumerate}
\end{env}

\subsection{Complétés de modules d'homomorphismes}
\label{subsection:III.4.5-fr}

\begin{proposition}[4.5.1]
\label{III.4.5.1-fr}
Soient $A$ un anneau noethérien, $\mathfrak J$ un idéal de $A$, $X$ un
$A$-préschéma de type fini, $\mathcal F$, $\mathcal G$ deux
$\mathcal O_X$-Modules cohérents dont les supports ont une intersection
propre sur $Y=\operatorname{Spec}(A)$ (II, 5.4.10). Alors, pour tout entier
$n\geq0$, $\operatorname{Ext}^n_{\mathcal O_X}(X;\mathcal F,\mathcal G)$ est un $A$-module
de type fini, et son séparé complété pour la topologie $\mathfrak J$-préadique
s'identifie canoniquement (avec les notations de
(\hyperref[III.4.1.7-fr]{4.1.7})) à
$\operatorname{Ext}^n_{\mathcal O_{\widehat X}}(\widehat X;\widehat{\mathcal F},
\widehat{\mathcal G})$.
\end{proposition}

\begin{proof}
On sait (T, 4.2) qu'il existe une suite spectrale birégulière
$\mathrm E(\mathcal F,\mathcal G)$ dont l'aboutissement est
$\operatorname{Ext}_{\mathcal O_X}(X;\mathcal F,\mathcal G)$ et dont les termes
$\mathrm E_2$ sont donnés par
\[
\mathrm E_2^{pq}=\mathrm H^p\bigl(X,
\mathcal{E}\!xt^q_{\mathcal O_X}(\mathcal F,\mathcal G)\bigr).
\]
On sait que $\mathcal{E}\!xt^q_{\mathcal O_X}(\mathcal F,\mathcal G)$ est un
$\mathcal O_X$-Module cohérent (0, 12.3.3) dont le support est contenu dans
l'intersection de ceux de $\mathcal F$ et $\mathcal G$ (T, 4.2.2), et est
par suite propre sur $Y$ (II, 5.4.10). On conclut de
(\hyperref[III.3.2.4-fr]{3.2.4}) que les $\mathrm E_2^{pq}$ sont des
$A$-modules de type fini, et par suite (0, 11.1.8) il en est de même de tous
les termes $\mathrm E_r^{pq}$ de la suite spectrale et de son aboutissement.
D'autre part, si $i:\widehat X\to X$ est le morphisme canonique,
$\widehat{\mathcal F}$ et $\widehat{\mathcal G}$ s'identifient canoniquement
à $i^*(\mathcal F)$ et $i^*(\mathcal G)$, et $i$ est plat (I, 10.8.8 et
10.8.9). On sait alors (0, 12.3.4) qu'il existe, pour tout $q\geq0$, un
$i$-morphisme canonique
\[
u_q:\mathcal{E}\!xt^q_{\mathcal O_X}(\mathcal F,\mathcal G)
\longrightarrow
\mathcal{E}\!xt^q_{\mathcal O_{\widehat X}}
(\widehat{\mathcal F},\widehat{\mathcal G}),
\]
et que le $\mathcal O_{\widehat X}$-homomorphisme correspondant
\[
u_q^\#:i^*\bigl(\mathcal{E}\!xt^q_{\mathcal O_X}(\mathcal F,\mathcal G)\bigr)
\longrightarrow
\mathcal{E}\!xt^q_{\mathcal O_{\widehat X}}
(\widehat{\mathcal F},\widehat{\mathcal G})
\]
est un isomorphisme (0, 12.3.5)~; autrement dit (I, 10.8.8),
$\mathcal{E}\!xt^q_{\mathcal O_{\widehat X}}
(\widehat{\mathcal F},\widehat{\mathcal G})$ s'identifie canoniquement au
complété $(\mathcal{E}\!xt^q_{\mathcal O_X}(\mathcal F,\mathcal G))^\wedge$ (avec les
notations de (\hyperref[III.4.1.7-fr]{4.1.7})). On conclut alors du th. de
comparaison (\hyperref[III.4.1.10-fr]{4.1.10}) que pour tout $p\geq0$,
\[
\mathrm H^p\bigl(\widehat X,
\mathcal{E}\!xt^q_{\mathcal O_{\widehat X}}
(\widehat{\mathcal F},\widehat{\mathcal G})\bigr)
\]
s'identifie canoniquement au séparé complété de
$\mathrm H^p\bigl(X,\mathcal{E}\!xt^q_{\mathcal O_X}
(\mathcal F,\mathcal G)\bigr)$ pour la topologie $\mathfrak J$-préadique.

